%%%%%%%%%%%%%%%%%%%%%%%%%%%%%%%%%%%%%%%%%%%%%%%%%%%%%%%%%%%%%%%%%%%%%%%%%%%
% Función: aspectos éticos, económicos, sociales y ambientales del trabajo.
%
% Uso: adapta los apartados a tu proyecto. El manual incluye este apéndice
% únicamente para titulaciones de la ETSIT-UPM.
%
% Fuente: anexo A de la plantilla de TFG de la ETSIT-UPM para GISTS y GIB.
%%%%%%%%%%%%%%%%%%%%%%%%%%%%%%%%%%%%%%%%%%%%%%%%%%%%%%%%%%%%%%%%%%%%%%%%%%%

\chapter{Aspectos éticos, económicos, sociales y ambientales}
\label{cha:aspectos-eticos}

El apartado \emph{Requisitos de las acreditaciones internacionales EUR-ACE y ABET} de la Normativa de TFT de la ETSIT-UPM establece que ``La memoria del TFT del GITST, GIB y MUIT, y en general la de aquellas titulaciones que hayan obtenido o para las que se desee solicitar una acreditación internacional EUR-ACE o ABET, debe mostrar conciencia de la responsabilidad de la aplicación práctica de la ingeniería, el impacto social y ambiental, el compromiso con la ética profesional, la responsabilidad y las normas de la aplicación práctica de la ingeniería, así como sobre las prácticas de gestión de proyectos, gestión, y control de riesgos, entendiendo sus limitaciones''.

Este anexo obligatorio del TFT tendrá un carácter sintético con los siguientes apartados:

\section{Introducción}
\label{sec:aspectos-eticos-introduccion}

Breve descripción del contexto del proyecto, objetivos, necesidades que pretende cubrir o problemas que pretende resolver, centrándose en su relación con los temas sociales, económicos, éticos, legales y/o ambientales que se hayan identificado.

\section{Descripción de impactos relevantes relacionados con el proyecto}
\label{sec:aspectos-eticos-impactos}

Síntesis del trabajo realizado en la fase 2, de selección y descripción de impactos. Presentar y justificar las conclusiones a las que se haya llegado sobre cuáles son los asuntos más relevantes relacionados con la sostenibilidad social, económica o ambiental, así como los principales grupos de interés identificados y que se han considerado en los análisis posteriores.

\section{Análisis detallado de alguno de los principales impactos}
\label{sec:aspectos-eticos-analisis}

Síntesis del trabajo de análisis realizado.

\section{Conclusiones}
\label{sec:aspectos-eticos-conclusiones}

Valorar el proyecto desde un punto de vista ético, social, económico y medioambiental y justificar si el uso de criterios de sostenibilidad ha aportado o puede aportar valor añadido al proyecto.

\section*{Lecturas complementarias para preparar tu análisis}

Las indicaciones anteriores reproducen la estructura de la plantilla ETSIT-UPM. Las siguientes lecturas son sugerencias nuestras para ayudarte a desarrollarla, no requisitos adicionales de esa normativa. Puedes apoyarte en el código de conducta de ALLEA~\cite{allea2023integrity} para reflexionar sobre la integridad de la investigación y en la Agenda 2030~\cite{un2030agenda} para situar los posibles impactos sociales y ambientales de tu proyecto.

Como ejemplo de impacto social, los trabajos sobre traducción de habla a lengua de signos~\cite{sansegundo2008speech,sansegundo2008proposing} te permiten plantear preguntas sobre accesibilidad, usuarios destinatarios y límites de las soluciones técnicas. No se trata de que tu proyecto aborde ese mismo problema, sino de que justifiques a quién beneficia, qué dificultades puede introducir y cómo vas a valorar esos efectos.

%%% Local Variables:
%%% TeX-master: "../book"
%%% End:
